\appendix
\section{附录：程序代码}

本附录给出建模、特征识别、潜在表型分析与权重敏感性分析所用的完整代码。XGBoost 补充二分类模型位于代码二，XGBoost 特征重要性与 SHAP 解释性分析位于代码三。为避免正文被程序实现细节打断，代码统一置于附录。

\subsection{代码一：数据预处理与质量控制}
\VerbatimInput[fontsize=\tiny,breaklines=true,breakanywhere=true,numbers=left,numbersep=4pt,xleftmargin=1.8em]{../src/stage2_qc_pipeline.py}

\clearpage
\subsection{代码二：二分类辅助诊断模型与 XGBoost 补充模型}
\VerbatimInput[fontsize=\tiny,breaklines=true,breakanywhere=true,numbers=left,numbersep=4pt,xleftmargin=1.8em]{../src/stage3_baseline_models.py}

\clearpage
\subsection{代码三：判别性声学特征识别与 XGBoost/SHAP 解释性分析}
\VerbatimInput[fontsize=\tiny,breaklines=true,breakanywhere=true,numbers=left,numbersep=4pt,xleftmargin=1.8em]{../src/stage4_biomarker_identification.py}

\clearpage
\subsection{代码四：二类潜在语音表型建模}
\VerbatimInput[fontsize=\tiny,breaklines=true,breakanywhere=true,numbers=left,numbersep=4pt,xleftmargin=1.8em]{../src/stage5_pd_two_cluster_phenotyping.py}

\clearpage
\subsection{代码五：六类潜在语音表型建模}
\VerbatimInput[fontsize=\tiny,breaklines=true,breakanywhere=true,numbers=left,numbersep=4pt,xleftmargin=1.8em]{../src/stage6_pd_six_cluster_phenotyping.py}

\clearpage
\subsection{代码六：PD 内部分型特征扩展与权重敏感性分析}
\VerbatimInput[fontsize=\tiny,breaklines=true,breakanywhere=true,numbers=left,numbersep=4pt,xleftmargin=1.8em]{../src/stage5b_stage6b_pdonly_sensitivity.py}
